%File: main.tex
\documentclass[letterpaper]{article} % DO NOT CHANGE THIS
\usepackage{aaai2026}
\usepackage{times}  % DO NOT CHANGE THIS
\usepackage{helvet}  % DO NOT CHANGE THIS
\usepackage{courier}  % DO NOT CHANGE THIS
\usepackage[hyphens]{url}  % DO NOT CHANGE THIS
\usepackage{graphicx} % DO NOT CHANGE THIS
\urlstyle{rm} % DO NOT CHANGE THIS
\def\UrlFont{\rm}  % DO NOT CHANGE THIS
\usepackage{natbib}  % DO NOT CHANGE THIS AND DO NOT ADD ANY OPTIONS TO IT
\usepackage{caption} % DO NOT CHANGE THIS AND DO NOT ADD ANY OPTIONS TO IT
\frenchspacing  % DO NOT CHANGE THIS
\setlength{\pdfpagewidth}{8.5in} % DO NOT CHANGE THIS
\setlength{\pdfpageheight}{11in} % DO NOT CHANGE THIS
%
% These are recommended to typeset algorithms but not required. See the subsubsection on algorithms. Remove them if you don't have algorithms in your paper.
\usepackage{algorithm}
\usepackage{algorithmic}

%
% These are are recommended to typeset listings but not required. See the subsubsection on listing. Remove this block if you don't have listings in your paper.
\usepackage{newfloat}
\usepackage{listings}
\DeclareCaptionStyle{ruled}{labelfont=normalfont,labelsep=colon,strut=off} % DO NOT CHANGE THIS
\lstset{%
	basicstyle={\footnotesize\ttfamily},% footnotesize acceptable for monospace
	numbers=left,numberstyle=\footnotesize,xleftmargin=2em,% show line numbers, remove this entire line if you don't want the numbers.
	aboveskip=0pt,belowskip=0pt,%
	showstringspaces=false,tabsize=2,breaklines=true}
\floatstyle{ruled}
\newfloat{listing}{tb}{lst}{}
\floatname{listing}{Listing}
%
% Keep the \pdfinfo as shown here. There's no need
% for you to add the /Title and /Author tags.
\pdfinfo{
/TemplateVersion (2026.1)
}

% DISALLOWED PACKAGES
% \usepackage{authblk} -- This package is specifically forbidden
% \usepackage{balance} -- This package is specifically forbidden
% \usepackage{color (if used in text)
% \usepackage{CJK} -- This package is specifically forbidden
% \usepackage{float} -- This package is specifically forbidden
% \usepackage{flushend} -- This package is specifically forbidden
% \usepackage{fontenc} -- This package is specifically forbidden
% \usepackage{fullpage} -- This package is specifically forbidden
% \usepackage{geometry} -- This package is specifically forbidden
% \usepackage{grffile} -- This package is specifically forbidden
% \usepackage{hyperref} -- This package is specifically forbidden
% \usepackage{navigator} -- This package is specifically forbidden
% (or any other package that embeds links such as navigator or hyperref)
% \indentfirst} -- This package is specifically forbidden
% \layout} -- This package is specifically forbidden
% \multicol} -- This package is specifically forbidden
% \nameref} -- This package is specifically forbidden
% \usepackage{savetrees} -- This package is specifically forbidden
% \usepackage{setspace} -- This package is specifically forbidden
% \usepackage{stfloats} -- This package is specifically forbidden
% \usepackage{tabu} -- This package is specifically forbidden
% \usepackage{titlesec} -- This package is specifically forbidden
% \usepackage{tocbibind} -- This package is specifically forbidden
% \usepackage{ulem} -- This package is specifically forbidden
% \usepackage{wrapfig} -- This package is specifically forbidden
% DISALLOWED COMMANDS
% \nocopyright -- Your paper will not be published if you use this command
% \addtolength -- This command may not be used
% \balance -- This command may not be used
% \baselinestretch -- Your paper will not be published if you use this command
% \clearpage -- No page breaks of any kind may be used for the final version of your paper
% \columnsep -- This command may not be used
% \newpage -- No page breaks of any kind may be used for the final version of your paper
% \pagebreak -- No page breaks of any kind may be used for the final version of your paperr
% \pagestyle -- This command may not be used
% \tiny -- This is not an acceptable font size.
% \vspace{- -- No negative value may be used in proximity of a caption, figure, table, section, subsection, subsubsection, or reference
% \vskip{- -- No negative value may be used to alter spacing above or below a caption, figure, table, section, subsection, subsubsection, or reference

\setcounter{secnumdepth}{0} %May be changed to 1 or 2 if section numbers are desired.

% The file aaai2026.sty is the style file for AAAI Press
% proceedings, working notes, and technical reports.
%

% Title

% Your title must be in mixed case, not sentence case.
% That means all verbs (including short verbs like be, is, using,and go),
% nouns, adverbs, adjectives should be capitalized, including both words in hyphenated terms, while
% articles, conjunctions, and prepositions are lower case unless they
% directly follow a colon or long dash
\title{Security Threats and AI Defenses in 6G\\\LaTeX{}}
\author{
    %Authors
    % All authors must be in the same font size and format.
    Written by: \\
    Jodhveer Atwal,
    Khushi Malik,
    Mary Klassen,
    Arda Yazici,
    Victor Mai
}
\date{
    \textsuperscript{\rm 1}
}

%Example, Single Author, ->> remove \iffalse,\fi and place them surrounding AAAI title to use it
\iffalse
\title{My Publication Title --- Single Author}
\author {
    Author Name
}
\date{
    Kwantlen Polytechnic University\\
    Surrey, British Columbia, Canada\\
}
\fi

\iffalse
%Example, Multiple Authors, ->> remove \iffalse,\fi and place them surrounding AAAI title to use it
\title{My Publication Title --- Multiple Authors}
\author {
    % Authors
    First Author Name\textsuperscript{\rm 1},
    Second Author Name\textsuperscript{\rm 2},
    Third Author Name\textsuperscript{\rm 1}
}
\affiliations {
    % Affiliations
    \textsuperscript{\rm 1}Affiliation 1\\
    \textsuperscript{\rm 2}Affiliation 2\\
    firstAuthor@affiliation1.com, secondAuthor@affilation2.com, thirdAuthor@affiliation1.com
}
\fi


% REMOVE THIS: bibentry
% This is only needed to show inline citations in the guidelines document. You should not need it and can safely delete it.
\usepackage{bibentry}
% END REMOVE bibentry

\begin{document}

\maketitle

\begin{abstract}
This paper surveys security threats and AI-based defense strategies in emerging 6G networks. As 6G systems increasingly rely on artificial intelligence for autonomous decision-making, resource management, and network optimization, they also introduce new attack surfaces such as adversarial evasion, poisoning, privacy leakage, and model manipulation.
\end{abstract}

\section{Introduction}

The transition from 5G to 6G represents a fundamental shift from human-managed connectivity to AI-native networking. To support the extreme requirements of 6G such as terahertz (THz) frequencies and massive-scale IoT networks must become "Zero-Touch", utilizing artificial intelligence (AI) to autonomously handle resource allocation and network management. However, this heavy reliance on AI introduces a novel and dangerous attack surface where the AI itself becomes the primary target.

This report explores the security landscape of AI-enabled 6G networks through a multi-layered lens. We begin by establishing the technical background of 6G and O-RAN, specifically focusing on how adversarial machine learning (AML) including data poisoning and evasion attacks---can disrupt intelligent wireless networking. We then provide a comprehensive literature review covering a broad spectrum of adversarial risks, from specific wireless attack scenarios to the evolving trends in countermeasures.

A significant portion of our analysis is dedicated to emerging defense frameworks. We investigate the transition toward a Zero Trust Architecture (ZTA), where "implicit trust" is replaced by continuous AI-based verification of identity and communication. Furthermore, we discuss the implementation of Explainable AI (XAI) and layered defenses to secure the 6G-IoT ecosystem. By comparing current defense strategies and identifying existidng research gaps, this paper aims to define a path toward a resilient, self-healing, and secure 6G infrastructure.

\section{Background}

%%% -----------------------------------------------------
%%% SECTION: Background
%%% AUTHOR : Jodhveer (@Jatwal27)
%%% -----------------------------------------------------

\subsection{The Evolution Toward 6G and O-RAN}

The shift from 5G to 6G represents a significant evolution in how global connectivity is designed. While 5G improved throughput, latency, and network virtualization, 6G aims to extend connectivity across integrated terrestrial, aerial, and satellite networks to enable three-dimensional coverage \cite{kaur2025}. By linking ground infrastructure with non-terrestrial systems, 6G seeks to reduce connectivity gaps across geographic regions. This architecture allows devices to dynamically transition between terrestrial towers and satellite systems. Research visions for 6G propose peak data rates approaching 1 Tbps and support for up to $10^6$ IoT devices per $\text{km}^2$ \cite{kaur2025}.

To support networks of this scale, the telecommunications industry is moving toward Open Radio Access Network (O-RAN) architectures. Traditionally, network hardware and software were tightly coupled and typically supplied by a single vendor. O-RAN introduces standardized interfaces that allow operators to deploy interoperable components from multiple vendors, increasing flexibility while also introducing new security considerations \cite{kaur2025}. Within this architecture, artificial intelligence assists network management through the RAN Intelligent Controller (RIC). Near-real-time applications called xApps handle short-term operational tasks such as signal optimization, while rApps support longer-term network management and policy decisions.

\subsection{AML Threats in AI-Driven Networks}

Because artificial intelligence plays a central role in managing 6G networks, the machine learning (ML) models themselves become potential targets for attackers \cite{Catak_Catak_Moldsvor_2021}. Adversarial Machine Learning (AML) studies how malicious actors manipulate ML systems to disrupt automated decision making.

One common technique is the evasion attack, where attackers introduce small disruptions into input data during inference to cause incorrect predictions. These disruptions exploit the sensitivity of ML models and can lead to incorrect outputs without altering the training data \cite{Catak_Catak_Moldsvor_2021}. In network environments, this may result in misclassification of traffic patterns or incorrect resource allocation.

Another major threat is data poisoning, in which attackers manipulate the training data used to build ML models. By injecting malicious samples, attackers can alter the learned decision boundaries and hurt model performance \cite{Chen_Zhang_Zhang_Wang_Liu_2021}. A well-known example is the clean-label poisoning attack, where malicious samples appear legitimate to human reviewers but introduce hidden patterns that trigger incorrect model behavior after deployment \cite{NEURIPS2018_22722a34}.

\subsection{Zero Trust Security and Explainable AI in 6G Networks}

The open and software-driven design of 6G networks fundamentally changes how security is implemented. Traditional network security relied on a perimeter-based approach in which internal systems were assumed to be trusted. However, 6G networks are highly distributed across cloud infrastructure, edge devices, satellites, and billions of connected devices, making a single security boundary impractical \cite{kaur2025}. 

To address this challenge, many 6G security models adopt Zero Trust Architecture (ZTA), where no device or user is trusted by default. Every request must be continuously verified based on identity, behavior, and context \cite{ali2025, kaur2025}. Given the scale and complexity of 6G environments, AI-driven monitoring is essential for analyzing traffic patterns and detecting anomalies in real time.

At the same time, the increasing reliance on machine learning introduces the need for transparency in automated decisions. Many ML models operate as complex ``black boxes,'' making it difficult to understand why specific actions are taken. This lack of transparency can weaken trust and hinder effective security management. Explainable Artificial Intelligence (XAI) addresses this issue by providing interpretable insights into model behavior \cite{kaur2025}.

Common XAI techniques include SHAP, which quantifies the contribution of each feature to a prediction, and LIME, which explains individual predictions using local approximations of the model. These methods help administrators audit AI decisions, detect abnormal behavior, and improve trust in AI-driven network security systems \cite{kaur2025}.


\subsection{Zero Trust Security and Explainable AI in 6G Networks}

The open and software-driven design of 6G networks fundamentally changes how security is implemented. Traditional network security relied on a perimeter-based approach in which internal systems were assumed to be trusted. However, 6G networks are highly distributed across cloud infrastructure, edge devices, satellites, and billions of connected devices, making a single security boundary impractical \cite{kaur2025} .
To address this challenge, many 6G security models adopt Zero Trust Architecture (ZTA), where no device or user is trusted by default. Every request must be continuously verified based on identity, behaviour, and context \cite{kaur2025} . Given the scale and complexity of 6G environments, AI-driven monitoring is essential for analyzing traffic patterns and detecting anomalies in real time.

At the same time, the increasing reliance on machine learning introduces the need for transparency in automated decisions. Many ML models operate as complex “black boxes,” making it difficult to understand why specific actions are taken. This lack of transparency can weaken trust and hinder effective security management. Explainable Artificial Intelligence (XAI) addresses this issue by providing interpretable insights into model behaviour \cite{kaur2025} .

Common XAI techniques include SHAP, which quantifies the contribution of each feature to a prediction, and LIME, which explains individual predictions using local approximations of the model. These methods help administrators audit AI decisions, detect abnormal behaviour, and improve trust in AI-driven network security systems \cite{kaur2025}.

\section{Literature Review}

%%% -----------------------------------------------------
%%% SECTION: Literature Review
%%% AUTHOR : Arda (@snowcactus)
%%% -----------------------------------------------------

\subsection{Zero Trust, O-RAN Security, and AI-Based Communication}

The transition to 6G introduces a degree of openness and heterogeneity that renders traditional "castle-and-moat" security perimeters obsolete. As established in recent studies on Software-Defined Zero Trust Architecture (ZTA), the 6G ecosystem must move toward a decentralized access control regime. This model treats network control domains as "communities," using adaptive collaboration to mitigate large-scale threats such as Distributed Denial of Service (DDoS) attacks and zero-day exploits that leverage the increased use of open-source software in 6G network entities.

A critical vulnerability exists within the Open Radio Access Network (O-RAN) architecture. While O-RAN promotes vendor interoperability, it also introduces new forward interfaces (like the E2 and A1 interfaces) that are susceptible to adversarial manipulation. Research into O-RAN security for 6G highlights that the Near-Real-Time RAN Intelligent Controller (Near-RT RIC) is a high-risk target. Attackers can potentially exploit the xApps and rApps that manage network optimization. To counter this, O-RAN standards are evolving to include PKI-based mTLS for the management plane and AI-driven anomaly detection to secure internal API communications.

Furthermore, the integration of 6G-IoT networks requires a shift toward AI-based Intrusion Detection Systems (IDS) that prioritize privacy. Current proposals suggest using an AIBID-SCSA (AI-Based Intrusion Detection and Secure Communication) technique, which employs algorithms like the Improved Sparrow Search Algorithm (ISSA) to format and recognize attack patterns across dense IoT clusters. By leveraging Federated Learning (FL), these systems can train on distributed edge nodes, allowing the network to recognize dynamic cyber-attacks without exposing raw user data to the central cloud. This ensures that 6G communication remains "zero-touch" while maintaining strict data confidentiality.

\subsection{Adversarial Machine Learning and Data Poisoning Threats \& Proactive AI-Driven Mitigations}

A primary concern for 6G security is the vulnerability of the underlying Machine Learning (ML) models to data poisoning. Unlike traditional cyberattacks that target software bugs, poisoning attacks target the "logic" of the network. As demonstrated in the "Poison Frogs!" study, clean-label poisoning allows attackers to inject malicious samples into training sets without altering the labels. This is particularly dangerous for 6G because the data remains visually and logically consistent, making it nearly impossible for manual audits to detect. In a 6G context, an attacker could subtly feed altered data into a network’s training set to force the model to misclassify specific traffic as "safe" during the operational phase.

The impact of these attacks extends to the physical layer of the network. Research into mmWave Beam Prediction shows that adversarial attacks, such as the Fast Gradient Sign Method (FGSM), can manipulate radio frequency data. By adding a small amount of calculated noise to a signal, an attacker can trick a 6G base station into misdirecting its beams. This results in dropped connections and severe latency, proving that adversarial ML is not merely a software vulnerability but a direct threat to 6G physical connectivity.


To counter these evolving threats, the literature suggests a shift toward attack-agnostic defense systems. The De-Pois framework represents a significant advancement in this area by utilizing Generative Adversarial Networks (GANs) to establish a "clean" baseline of data distribution. By comparing real-time training samples against this GAN-generated baseline, the system can filter out malicious data without needing prior knowledge of the specific attack type. This is a crucial defense for 6G’s autonomous policy-making models, ensuring that routing and security protocols remain uncompromised.

However, the integration of these defenses is not without challenges. While AI-driven security is necessary to handle the complexity of 6G resource management, it introduces a "double-edged sword" scenario. Advanced models like De-Pois or layered security frameworks require significant computational power, which could introduce unwanted latency or exhaust the battery life of small IoT devices. Furthermore, many of these AI-based security measures remain in the early research phase and have yet to be tested against the high-speed, dynamic environment of a real-world 6G deployment.



\section{Literature Review}

%%% -----------------------------------------------------
%%% SECTION: Literature Review
%%% AUTHOR : Khushi (@kelsiei)
%%% -----------------------------------------------------

\subsection{Adversarial ML Threats and Defenses}

Recent research shows that artificial intelligence is becoming central to 6G because it helps future wireless networks adapt more quickly, automate decision-making, and manage increasingly complex environments \cite{bountakas2023defense,xu2022security}. At the same time, this shift expands the security problem: attackers can now target not only network infrastructure, but also the machine-learning models that guide network behavior. Broad adversarial machine learning surveys show that evasion and poisoning attacks remain major concerns, while 6G-focused studies argue that AI-enabled optimization introduces new risks in areas such as resource allocation, autonomous control, and distributed intelligence \cite{bountakas2023defense,xu2022security}. As a result, security in 6G must be studied as both a networking problem and a learning-system problem.

This concern becomes more concrete in 6G-specific architectures that rely on distributed intelligence and open control frameworks. Xu, Su, and Li describe AI-enabled 6G using a space-air-ground-ocean integrated network (SAGOIN) architecture and two federated learning (FL) training models, namely single-region FL and cross-region FL, to support large-scale intelligent services while reducing raw-data sharing \cite{xu2022security}. However, they also show that this design introduces three major threats: malicious cheating attacks, low-quality local model training attacks, and privacy stealing attacks \cite{xu2022security}. Similar concerns appear in O-RAN-based cellular systems, where machine-learning xApps inside the near-real-time RAN Intelligent Controller can be manipulated through adversarial inputs, causing sharp drops in model accuracy unless robustness measures such as distillation are applied \cite{chiejina2024oran}. Together, these studies suggest that 6G defenses must be domain-aware, fast enough for operational deployment, and robust against attacks aimed directly at learning pipelines.

Although existing work identifies important threats, the defense landscape is still fragmented. General adversarial machine learning surveys describe broad categories of proactive and reactive defenses, but these methods are often evaluated outside realistic wireless environments \cite{bountakas2023defense}. In contrast, wireless-focused studies propose more domain-specific protections for signal classifiers and AI-driven radio intelligence, yet they are usually tested in narrower settings \cite{ahmed2023cdi,chiejina2024oran}. Xu, Su, and Li partially bridge this gap by proposing a defense framework that combines trust evaluation, Q-learning-based incentives, and local differential privacy, but their solution is still presented as a case study rather than a general deployment-ready framework \cite{xu2022security}. This creates a gap between theory and deployment, because a defense that performs well in a controlled experiment may not remain effective under dynamic channel conditions, open-radio architectures, or latency-sensitive decision pipelines.


Overall, the literature suggests that security for AI-driven 6G remains an open problem rather than a solved one. Current studies provide useful building blocks for robustness, poisoning resistance, and trust-aware design, but they do not yet offer a unified defense framework that is scalable, explainable, and practical for real-time 6G systems \cite{xu2022security,bountakas2023defense,chiejina2024oran,ahmed2023cdi}. This gap motivates the need for future work that combines adversarial robustness, secure distributed learning, and operationally realistic defense strategies.

%%% -----------------------------------------------------
%%% SECTION: Literature Review
%%% AUTHOR : Victor (@lugioo)
%%% -----------------------------------------------------

\subsection{Layered Defenses}

Since 6G networks are so complex, researchers have largely moved away from the idea that any single security technology can protect them fully. Kumar et al. (2026) argues that the right approach is to stack complementary mechanisms across different levels of the network so that weaknesses in one layered are covered by strengths in another. Their proposed framework organizes these mechanisms into three logical tiers: a foundational 6G infrastructure layer, an AI enabled security layer and an application and service layer, that has real time data flows with a AI powered decision loops connecting each tier. At the lowest level, physical layer security (PLS) exploits the natural properties of wireless channels, things like noise and signal fading which are used to protect transmissions without relying on traditional encryption. This makes it computionally cheap and naturally resistant to quantum attacks, but it cannot handle threats that originate higher up in the network like malware or poisoned training data. Moving up the stack, Distributed Ledger Technologies (DLT) addresses the need for trustworthy record keeping, providing temper evident logs of network activity that are particularly useful for identity management, though consensus latency makes it unsuitable for ultra low latency applications. Tying these layers together, distributed and federated AI/ML serves as what they call the unifying enabler across all other technologies, capable of detecting anomalies, predicting threats and orchestrating adaptive responses in real time. They back this up with evidence from the paper, pointing to hybrid CNN-LSTM intrusion detection models that achieve roughly 97-99\% detection accuracy on benchmark datasets, while reinforcement learning based can isolate compromised networks slices in sub-second time. To ground the whole framework into something practical, Kumar et al. walk through the examples of EV charging infrastructure in a smart city, showing how physical layer federated detection can all work together in a single deployment. The takeaway is that layered defense in 6G is less about picking the best technology and more about understanding how different tools complement each other across the network stack.


%%% -----------------------------------------------------
%%% SECTION : Literature Review
%%% AUTHOR  : Mary (@MustangKoschei)
%%% -----------------------------------------------------
\subsection{Adversarial Risks}

The transition from 5G to 6G intelligence significantly expands the network's attack surface as it moves towards AI-native, "Zero-Touch" networking,where AI is ultilized to autonomously handle critical functions. Now, security threats now target not only the physical infrastructure but also the machine-learning models governing network behavior. 
One prominent 6G-specific risk is adversarial beam manipulation within the physical layer. Research demonstrates that techniques like the Fast Gradient Sign Method (FGSM) can be used to inject calculated noise into radio frequency data sent to these models. This can deceive base stations into misdirecting signals, causing connections to drop or data rates to plummet without triggering traditional security alarms\cite{Catak_Catak_Moldsvor_2021}. Furthermore, 6G's reliance on training data makes it highly vulnerable to "clean-label" poisoning attacks, where attackers subtly modify the features of training samples while keeping their labels consistent(Wang et al. 2021). These attacks are exceptionally difficult to detect because the malicious data appears visually and logically correct, yet it forces the model to misclassify specific targets during operation. In network-assisted IoT environments, these adversarial tactics have been shown to be more effective and damaging than traditional jamming at equivalent power levels\cite{son2024}

\subsection{Countermeasure Trends}  

To address these threats, researchers are shifting away from singular security technologies in favour of Intelligent Zero Trust Architecture (i-ZTA) and layered defense frameworks. The i-ZTA model is specifically designed to be O-RAN compliant, utilizing the Near-Real-Time RAN Intelligent Controller (RIC) to host AI engines that perform real-time risk assessments for every individual access request\cite{Ramezanpour_2022}. This approach facilitates "zero-touch" management and behavior-based isolation, which is critical for autonomous "tactical bubbles" that operate without on-site administrators\cite{Suomalainen_2025}. Technical trends also include the development of the "De-Pois" system, which employs Generative Adversarial Networks (GANs) to create a clean data distribution baseline, allowing the network to filter out poisoned training samples without prior knowledge of the specific attack type\cite{Chen_Zhang_Zhang_Wang_Liu_2021}. Furthermore, innovative frameworks like the Quantum-Enhanced Blockchain Framework (QEBF) are being proposed to protect intrusion detection systems by integrating blockchain and quantum principles to detect anomalies and restore model integrity\cite{Rajesh_2025}. Ultimately, the literature points toward a strategic roadmap for 2035 that focuses on developing verifiable and quantum-safe AI models capable of managing complex 6G resources while maintaining high detection accuracy\cite{Kumar_Dutta_Vamsi_Sankararao_Varri_Puthal_2026}.

%%% END SECTION: Literature Review
% Para 3
% Para 4

%%% END SECTION: Literature Review

% Uncomment the following to link to your code, datasets, an extended version or similar.
% You must keep this block between (not within) the abstract and the main body of the paper.
% \begin{links}
%     \link{Code}{https://aaai.org/example/code}
%     \link{Datasets}{https://aaai.org/example/datasets}
%     \link{Extended version}{https://aaai.org/example/extended-version}
% \end{links}

\section{Discussion}

%%% -----------------------------------------------------
%%% SECTION: Discussion
%%% AUTHOR : Victor (@lugioo)
%%% -----------------------------------------------------

Across the reviewed literature, a consistent theme emerges: as 6G networks delegate critical functions to AI, the attack surface shifts from physical infrastructure to the machine learning models that govern it. Across all reviewed works, adversarial threats such as evasion attacks, clean-label data poisoning and adversarial beam manipulation emerge as defining security challenges of the 6G era. What sets these threats apart from traditional cyberattacks is their ability to evade detection entirely, as malicious inputs can appear fully legitimate to human auditors. As a result, conventional monitoring approaches appear insufficient.

A recurring theme in the defense literature is that reliance on any single protection mechanism leaves critical vulnerabilities unaddressed. Kumar et al. argue for layered architectures that organize defenses across physical, distributed ledger and AI-enabled tiers. This view is reinforced by other works: Xu et al.'s federated learning framework addresses data privacy but still requires trust evaluation and incentive mechanisms to handle malicious participants, while Chiejina et al.'s distillation based approach for O-RAN xApps protects inference pipelines but leaves training time threats unaddressed. Similarly, the De-Pois GAN framework and Zero Trust Architecture each address distinct parts of the problem. Taken together, the literature suggests that robust 6G security is not a matter of selecting a single best defense, but of understanding how complementary mechanisms cover one another's blind spots across the full stack.

However, significant gaps still remain. First, almost all proposed frameworks are evaluated in controlled or simulated environments and none have been validated against the dynamic, high speed conditions of real world deployment. The jump from benchmark accuracy, such as the 97 to 99 percent intrusion detection rates cited by Kumar et al., to operational performance under adversarial real time conditions is challenging and largely unaddressed in the current literature. Second, many advanced defenses, particularly GAN based systems and federated learning pipelines incorporating differential privacy, carry high computational overhead. This poses a direct conflict with 6G's ultra low latency requirements and the resource constraints of IoT edge devices operating at the network boundary. Third, while Explainable AI (XAI) techniques such as SHAP and LIME are proposed to improve transparency and trust in AI driven security decisions, their integration into real time network management remains limited throughout the reviewed works. XAI is treated more as a desirable design property than a concretely deployed and evaluated mechanism.

Overall, the field lacks a unified defense framework that is simultaneously scalable, explainable and deployable in real time. Current proposals address subsets of these requirements in isolation from one another. Bridging this gap, particularly through work that combines adversarial robustness, secure distributed learning and operational practicality represents, the most critical priority for future 6G security research.




\section{Proposed Idea}

%%% -----------------------------------------------------
%%% SECTION: Proposed Idea
%%% AUTHOR : Jodhveer (@Jatwal27)
%%% -----------------------------------------------------

6G and O-RAN architectures rely heavily on artificial intelligence to automatically manage network operations. While this improves efficiency, it also introduces security risks. Machine learning models can be manipulated through adversarial inputs or poisoned data, leading to incorrect decisions that may spread across the network \cite{Catak_Catak_Moldsvor_2021, Chen_Zhang_Zhang_Wang_Liu_2021}.

Existing defenses are largely perimeter-based and focus on detecting threats at the network entry point. However, this creates a single point of failure, as attacks that bypass initial detection can directly impact AI-driven decision making. Given the distributed and software-based nature of 6G systems, a lifecycle-based defense approach is needed to protect AI across all stages of operation rather than relying on a single checkpoint \cite{Chen_Zhang_Zhang_Wang_Liu_2021}.

\subsection {Three Phase Defense}

This section proposes a layered framework that secures 6G AI systems across three stages: data, language model, and decision making.

\subsection {Securing the Data}

The first phase focuses on protecting data before it reaches the model. Incoming data is passed through a filtering stage that acts as a clean room, where a lightweight validation mechanism checks for signs of poisoning or abnormal patterns. This can be implemented using anomaly detection or statistical validation techniques to identify suspicious inputs and remove them before they affect the model \cite{Chen_Zhang_Zhang_Wang_Liu_2021}.

Some attacks are more difficult to detect because the data appears normal. Clean-label poisoning demonstrates that correctly labeled samples can still manipulate model behaviour during inference \cite{NEURIPS2018_22722a34}. To address this, an additional validation step can be used where a secondary model compares predictions and identifies inconsistencies between clean and potentially poisoned data. This approach is similar to De-Pois, where a mimic model is used to detect poisoned samples based on differences in model behaviour \cite{Chen_Zhang_Zhang_Wang_Liu_2021}.

This phase is especially important in 6G environments due to the large volume of data generated by IoT devices and distributed network components.

\subsection {Securing the AI}

The second phase focuses on improving the robustness of the AI model itself. Even with filtering, some malicious inputs may bypass detection, so the model must be able to operate under adversarial conditions.

This is achieved through adversarial training, where the model is trained using both normal data and adversarially crafted inputs. By exposing the model to small disturbances during training, it learns to resist evasion attacks that attempt to exploit its sensitivity to input changes \cite{Catak_Catak_Moldsvor_2021}.

This stress-testing process helps ensure that the model maintains reliable performance during real-time network operation, even when faced with manipulated inputs.

\subsection {Securing the Decisions}

The final phase focuses on securing the decisions made by the AI system. In 6G networks, where there is no clear security perimeter, Zero Trust is used to continuously verify all users and devices. Each request is evaluated based on behaviour and contextual information rather than assumed trust.

At the same time, Explainable Artificial Intelligence is used to provide transparency into these decisions. Techniques such as SHAP and LIME generate human-readable explanations that show why a model made a specific decision \cite{Kumar_Dutta_Vamsi_Sankararao_Varri_Puthal_2026}.

This allows administrators to verify whether a decision is the result of a real threat or a possible system error, improving trust and overall system reliability.


%%% -----------------------------------------------------
%%% SECTION : Conclution
%%% AUTHOR  : Mary (@MustangKoschei)
%%% -----------------------------------------------------
\section{Conclusion}
The transition to 6G marks a fundamental shift from human-managed connectivity to AI-native, "Zero-Touch" networking, where artificial intelligence autonomously handles critical functions, like resource allocation and network management\cite{Suomalainen_2025}. However, as this report has established, this heavy reliance on AI creates a dangerous attack surface where the learning systems themselves become the primary targets for manipulation\cite{Kumar_Dutta_Vamsi_Sankararao_Varri_Puthal_2026}. 
The research highlights that 6G-specific architectures are vulnerable to AML tactics that can have physical consequences. These risks include adversarial beam manipulation in the physical layer, and "clean-label" poisoning attacks, which can subtly alter model logic while appearing legitimate to human reviewers. These threats are particularly notable in O-RAN environments, where the Near-Real-Time RIC and its associated xApps serve as high-risk targets for adversarial manipulation\cite{Catak_Catak_Moldsvor_2021}.
To address these risks, this paper proposes a lifecycle-based Three Phase Defense framework. This includes targeting the data, the model, and the decisions it produces. First by securing the data by filtering out poisoned inputs before they reach the model. Next, securing the AI by utilizing adversarial training and stress-testing to ensure models remain robust against evasion attacks during real-time operation. And finally, securing the decisions by using ZTA with Explainable AI techniques like SHAP and LIME to provide the transparency needed for continuous verification and auditing\cite{Kumar_Dutta_Vamsi_Sankararao_Varri_Puthal_2026}.
Significant challenges remain regarding the computational overhead and processing latency introduced by these complex AI defenses\cite{Chen_Zhang_Zhang_Wang_Liu_2021}. Many current solutions remain theoretical and have yet to be validated against the dynamic channel conditions of real-world 6G deployments. Addressing these gaps remains essential to realizing the vision of a resilient, self-healing, and secure 6G network.




%%% -----------------------------------------------------
%%% SECTION : TEMPLATE
%%% AUTHOR  : useless stuff beyond this, delete it
%%% -----------------------------------------------------


\bibliography{references}

\end{document}